
@@M0@@
Il
faut trouver @@M1@@ tel que:
@@M2@@
Comme @@M3@@ admet un inverse à gauche @@M4@@, on a nécessairement
@@M5@@
Il faut donc montrer que @@M6@@ vérifie les relations (3.7.1).

Pour montrer que @@M7@@, on part de la relation @@M8@@ et on lui applique @@M9@@. On obtient donc: @@M10@@ par définition de @@M11@@. Or d'après la proposition 3.6, @@M12@@; d'où la relation voulue.

Pour achever la démonstration, il faut montrer que
@@M13@@
Posons @@M14@@. On obtient: @@M15@@. D'autre part, on a par hypothèse: @@M16@@, d'où d'après ce qui précède @@M17@@, et d'après la proposition 3.4: @@M18@@; par suite, on a également @@M19@@. Or @@M20@@. En effet, si @@M21@@, @@M22@@; or @@M23@@ (prop. 2.7), et pour que @@M24@@, il faut qu'il existe @@M25@@ tel que @@M26@@ (VI,5.10); en appliquant @@M27@@ ceci entraîne @@M28@@, donc @@M29@@.

Supposons maintenant @@M30@@ noethérien et @@M31@@ régulier. Alors @@M32@@ est aussi régulier puisque c'est un fibré projectif sur @@M33@@, et pour tout point @@M34@@ de @@M35@@, l'anneau local @@M36@@ de @@M37@@ en @@M38@@ est un anneau régulier, car @@M39@@ est quotient de @@M40@@ par un élément régulier (voir EGA @@M41@@ 17.1.8).

